\chapter{Action Representation: VLA의 진짜 병목}

\section{장의 목표}

VLA에서 가장 과소평가되기 쉬운 부분은 행동 표현(action representation)이다. 비전과 언어 backbone이 아무리 강해도, 로봇 행동을 잘못 표현하면 정책은 실제 제어에서 불안정해진다. 2024--2026년 VLA 논문군에서 action tokenization, action chunking, trajectory autoregressive modeling, diffusion action decoder가 계속 등장한 이유는 명확하다. VLA의 병목은 단순히 ``무엇을 이해하는가''가 아니라 ``이해한 것을 어떤 제어 변수로 바꾸는가''에 있다.

이 장은 2장에서 정의한 action space를 더 깊게 다룬다. 목표는 다음 네 질문에 답하는 것이다.

\begin{enumerate}
  \item continuous action과 discrete action token은 각각 어떤 장단점이 있는가?
  \item action chunking은 왜 속도 문제와 안정성 문제를 동시에 건드리는가?
  \item diffusion 또는 flow 기반 행동 생성은 VLA에 왜 적합한가?
  \item 좋은 action representation 논문을 평가하려면 어떤 실험이 필요한가?
\end{enumerate}

\begin{figure}[t]
  \centering
  \begin{tikzpicture}[node distance=6mm]
    \node[vlablock] (reg) {Regression\\\(a_t\in\mathbb{R}^d\)};
    \node[vlablock, right=of reg] (tok) {Token\\\(k_t\in\{1,\ldots,K\}\)};
    \node[vlablock, right=of tok] (chunk) {Chunk\\\(a_{t:t+H}\)};
    \node[vlablock, below=of tok] (diff) {Diffusion/Flow\\trajectory distribution};
    \node[vlablock, right=of diff] (spline) {Spline/Basis\\low-dimensional curve};
    \draw[vlaarrow] (reg) -- node[above,font=\sffamily\scriptsize]{discretize} (tok);
    \draw[vlaarrow] (tok) -- node[above,font=\sffamily\scriptsize]{sequence} (chunk);
    \draw[vlaarrow] (chunk) -- (spline);
    \draw[vlaarrow] (diff) -- (chunk);
  \end{tikzpicture}
  \caption{Author-created schematic. Action representation map. Regression, token, chunk, diffusion/flow, spline은 서로 배타적 분류가 아니라 VLA가 행동을 압축하고 복원하는 서로 다른 interface이다.}
  \label{fig:action_representation_map}
\end{figure}

\section{행동 표현이 어려운 이유}

언어는 이산 token sequence이다. 이미지는 patch token sequence로 바꿀 수 있다. 그러나 로봇 행동은 대개 연속값이며, 시간적으로 부드러워야 하고, robot embodiment와 controller에 의존한다. 이 차이가 VLA의 핵심 mismatch이다.

언어 모델 관점에서는 다음 예측이 자연스럽다.

\begin{equation}
  p_\theta(w_{1:L} \mid I, x)
  =
  \prod_{\ell=1}^{L}
  p_\theta(w_\ell \mid w_{<\ell}, I, x),
\end{equation}

여기서 $w_\ell$은 단어 또는 subword token이다. 반면 로봇 정책은 다음과 같은 연속 제어 문제에 가깝다.

\begin{equation}
  \act_t \in \mathbb{R}^{d_a},
  \qquad
  \act_{t+1} - \act_t \ \mathrm{is\ usually\ small}.
\end{equation}

즉, 로봇 행동은 값의 정확도와 시간적 smoothness가 모두 중요하다. 이산 token 예측을 그대로 가져오면 LLM/VLM 구조와는 잘 맞지만 quantization error가 생긴다. 연속 회귀를 쓰면 제어 변수에는 직접 맞지만 multi-modal 행동을 표현하기 어렵다.

\section{연속 행동 회귀}

가장 직접적인 방식은 정책이 연속 행동을 회귀하는 것이다.

\begin{equation}
  \hat{\act}_t
  =
  \mu_\theta(\hist_t,\lang),
  \qquad
  \act_t \in \mathbb{R}^{d_a}.
\end{equation}

학습은 평균제곱오차로 할 수 있다.

\begin{equation}
  \mathcal{L}_{\mathrm{reg}}
  =
  \sum_{i,t}
  \left\|
  \act_t^{(i)}
  -
  \mu_\theta(\hist_t^{(i)},\lang^{(i)})
  \right\|_2^2.
\end{equation}

장점은 단순성과 속도이다. policy head가 한 번의 forward pass로 행동을 내므로 real-time control에 유리하다. 또한 end-effector delta pose, joint velocity, gripper command 같은 기존 controller interface와 바로 연결된다.

약점은 multi-modality이다. 같은 관측과 같은 명령에서도 여러 행동이 모두 타당할 수 있다. 컵을 왼쪽으로 돌아 잡을 수도 있고 오른쪽으로 돌아 잡을 수도 있다. MSE loss는 그 평균을 선호하므로 물리적으로 어색한 중간 행동이 나올 수 있다. 이를 평균 행동 문제라고 부를 수 있다.

\begin{warningbox}{평균 행동 문제}
MSE 기반 연속 회귀는 demonstration distribution이 multi-modal일 때 각 mode 사이의 평균을 예측할 수 있다. 로봇에서는 이 평균이 충돌, 미끄러짐, 불안정한 grasp로 이어질 수 있다.
\end{warningbox}

\section{이산 action token}

이산 action token 방식은 연속 행동을 codebook index로 바꾼다.

\begin{equation}
  u_t = Q(\act_t),
  \qquad
  u_t \in \{1,\ldots,K\}.
\end{equation}

정책은 action token을 autoregressive하게 예측한다.

\begin{equation}
  p_\theta(u_{1:L}\mid \hist_t,\lang)
  =
  \prod_{\ell=1}^{L}
  p_\theta(u_\ell \mid u_{<\ell},\hist_t,\lang).
\end{equation}

이 방식의 매력은 VLM/LLM 구조와 잘 맞는다는 것이다. 언어 token, visual token, action token을 같은 sequence modeling 문제로 만들 수 있다. VQ-VLA류 논문은 이 지점에서 출발한다. 행동을 token으로 바꾸면 pretrained language model의 autoregressive machinery를 활용하기 쉽다.

그러나 action token에는 세 가지 문제가 있다. 첫째, codebook이 실제 로봇 행동의 geometry를 잘 보존해야 한다. 둘째, token sequence가 길어지면 inference latency가 증가한다. 셋째, token prediction error가 continuous control error로 어떻게 전파되는지 분석해야 한다.

Quantization error는 다음처럼 정의할 수 있다.

\begin{equation}
  \epsilon_Q
  =
  \mathbb{E}
  \left[
  \left\|
  \act - Q^{-1}(Q(\act))
  \right\|_2^2
  \right].
\end{equation}

하지만 $\epsilon_Q$가 낮다고 항상 좋은 것은 아니다. 중요한 것은 task-relevant error이다. 예를 들어 gripper command의 작은 오류는 치명적일 수 있지만, 이동 중 회전 각도의 작은 차이는 큰 문제가 아닐 수 있다. 따라서 action tokenizer는 단순 reconstruction error뿐 아니라 closed-loop success로 평가해야 한다.

이를 engineering diagnostic으로 쓰면 task-weighted action error를 다음처럼 볼 수 있다.

\begin{equation}
  \epsilon_{\mathrm{task}}
  =
  \mathbb{E}_{a\sim D}
  \left[
    (a-\hat{a})^\top
    W_{\mathrm{task}}
    (a-\hat{a})
  \right].
\end{equation}

여기서 \(W_{\mathrm{task}}\)는 실제 논문 objective가 아니라 action error의 task relevance를 설명하기 위한 conceptual weight이다. Contact phase, gripper timing, end-effector height error처럼 실패 비용이 큰 차원에는 더 큰 weight가 들어간다.

\section{Action chunking}

Action chunking은 한 번에 여러 step의 행동을 생성하는 방식이다.

\begin{equation}
  \Act_t
  =
  (\act_t,\act_{t+1},\ldots,\act_{t+H-1})
  \sim
  \pi_\theta(\cdot \mid \hist_t,\lang).
\end{equation}

이 방식은 VLA에서 특히 중요하다. 큰 VLA 모델은 매 제어 step마다 호출하기 어렵다. $H$ step chunk를 내면 모델 호출 빈도를 줄일 수 있다. controller frequency가 $r$ Hz이고 chunk length가 $H$이면 모델 호출 frequency는 대략 $r/H$가 된다.

\begin{equation}
  f_{\mathrm{model}}
  \approx
  \frac{r}{H}.
\end{equation}

그러나 $H$를 크게 잡으면 환경 변화에 반응이 늦어진다. 따라서 chunking은 latency와 reactivity의 trade-off이다.

\begin{equation}
  J(H)
  =
  \mathrm{SR}(H)
  -
  \lambda_1 \tau_{\mathrm{model}}(H)
  -
  \lambda_2 \mathrm{Staleness}(H).
\end{equation}

$\mathrm{Staleness}(H)$는 chunk가 길어질수록 현재 관측과 실행 중 action이 어긋나는 정도를 나타낸다. 좋은 action chunking 논문은 $H$를 고정하지 않고, task 난이도와 환경 변화 속도에 따른 민감도를 보여야 한다.

\section{Trajectory autoregressive modeling}

Trajectory autoregressive modeling은 단일 행동보다 긴 trajectory를 순차적으로 생성한다. action chunking과 비슷하지만, chunk 내부의 구조를 sequence model로 더 명시적으로 다룬다.

\begin{equation}
  p_\theta(\Act_t \mid \hist_t,\lang)
  =
  \prod_{h=0}^{H-1}
  p_\theta(\act_{t+h}
  \mid
  \act_{t:t+h-1}, \hist_t,\lang).
\end{equation}

이 방식은 행동 사이의 시간적 의존성을 표현할 수 있다. 예를 들어 grasp 접근, gripper close, lift, place는 순서가 중요하다. 단일 step 정책은 이 순서를 암묵적으로만 학습하지만 trajectory model은 chunk 내부 순서를 직접 모델링한다.

문제는 compounding error이다. 초반 행동이 틀리면 뒤 행동 조건도 틀어진다. 따라서 trajectory autoregressive 모델은 teacher forcing과 closed-loop rollout 사이의 gap을 확인해야 한다.

\begin{equation}
  \Delta_{\mathrm{rollout}}
  =
  \mathbb{E}[\ell_{\mathrm{closed}}]
  -
  \mathbb{E}[\ell_{\mathrm{teacher}}].
\end{equation}

$\Delta_{\mathrm{rollout}}$가 크면 offline training loss가 낮아도 실제 실행에서 무너질 수 있다.

\section{Diffusion action decoder}

Diffusion 방식은 행동 trajectory를 확률적으로 생성한다. 초기 noise trajectory $\Act^K$에서 시작해 denoising 과정을 거쳐 실행할 trajectory $\Act^0$를 얻는다.

\begin{equation}
  \Act^{k-1}
  =
  D_\theta(\Act^k, k, c_t),
  \qquad
  k=K,\ldots,1.
\end{equation}

여기서 $c_t$는 시각, 언어, 로봇 상태를 결합한 context이다. Diffusion action decoder의 장점은 multi-modal trajectory distribution을 표현하기 쉽다는 점이다. 같은 목표를 여러 방식으로 달성할 수 있을 때, 하나의 평균 trajectory가 아니라 가능한 trajectory sample을 생성할 수 있다.

단점은 inference cost이다. $K$번 denoising하면 모델 호출이 많아진다. 따라서 VLA에서 diffusion을 쓰려면 step 수를 줄이거나, flow matching, consistency model, distillation 같은 방법으로 빠르게 만들어야 한다. 2025--2026년 flow-based policy와 real-time diffusion policy가 중요해진 이유가 여기에 있다.

\section{B-spline과 저차원 행동 기저}

행동 trajectory를 모든 시점의 action으로 표현하면 차원이 커진다. B-spline 같은 기저 함수를 쓰면 trajectory를 몇 개의 control point로 표현할 수 있다.

\begin{equation}
  \act(t)
  =
  \sum_{j=1}^{M}
  \mathbf{c}_j B_j(t),
\end{equation}

여기서 $\mathbf{c}_j$는 control point이고 $B_j(t)$는 basis function이다. 이 방식은 행동을 부드럽게 만들고 token 수를 줄일 수 있다. BEAST류 접근은 이런 저차원 trajectory encoding의 장점을 활용한다.

그러나 spline 표현은 모든 과제에 맞지는 않는다. 접촉이 급격히 바뀌는 과제에서는 너무 부드러운 trajectory가 오히려 불리할 수 있다. 예를 들어 물체를 밀거나 삽입하는 과제에서는 순간적인 force change가 필요할 수 있다. 따라서 trajectory basis는 task dynamics와 함께 평가해야 한다.

\section{행동 표현과 embodiment}

행동 표현은 robot embodiment와 분리할 수 없다. 같은 end-effector delta pose라도 Franka arm, UR5, xArm에서 controller와 workspace가 다르다. Humanoid나 mobile manipulator로 가면 action space는 더 복잡해진다. Multi-embodiment VLA의 핵심 문제는 서로 다른 로봇의 action을 공통 표현으로 바꾸는 것이다.

이를 추상적으로 쓰면 각 robot $r$에 대해 action mapping이 있다.

\begin{equation}
  \act_t^{(r)}
  =
  M_r(z_t),
\end{equation}

여기서 $z_t$는 embodiment-independent latent action이고 $M_r$은 robot-specific decoder이다. 좋은 multi-embodiment VLA는 $z_t$가 task semantics를 담고, $M_r$이 embodiment constraint를 처리하도록 설계해야 한다. 그렇지 않으면 단순히 여러 로봇 데이터를 섞은 policy가 되어 일반화가 약해질 수 있다.

\section{평가 기준}

Action representation 논문은 다음 실험을 갖추어야 한다.

\begin{enumerate}
  \item 같은 backbone과 같은 데이터에서 action representation만 바꾼 비교
  \item reconstruction error와 task success의 관계 분석
  \item chunk length, codebook size, trajectory horizon 민감도
  \item inference latency와 memory 사용량
  \item closed-loop rollout에서의 failure mode
  \item unseen task와 unseen object에서의 generalization
\end{enumerate}

특히 reconstruction error만으로는 부족하다. 행동 표현이 좋아졌다는 주장은 실제 task success와 안정성으로 닫혀야 한다.

\begin{table}[t]
  \centering
  \caption{행동 표현 방식의 비교}
  \label{tab:action_rep_comparison}
  \begin{tabularx}{\textwidth}{>{\bfseries}p{30mm}XXX}
    \toprule
    표현 & 장점 & 약점 & 확인할 실험 \\
    \midrule
    Continuous regression & 빠르고 controller와 직접 연결됨 & multi-modal 행동에 약함 & MSE와 closed-loop success 비교 \\
    Action token & VLM/LLM 구조와 잘 맞음 & quantization error와 긴 sequence 문제 & codebook size, reconstruction, success \\
    Action chunk & 모델 호출 빈도 감소 & 반응성 저하 가능 & chunk length ablation \\
    Diffusion trajectory & multi-modal trajectory 표현력 & denoising latency & step 수와 success/latency trade-off \\
    Spline basis & 부드럽고 compact한 표현 & 급격한 접촉 변화에 약할 수 있음 & contact-rich task 평가 \\
    \bottomrule
  \end{tabularx}
\end{table}

\section{요약}

VLA의 행동 표현은 단순 구현 세부사항이 아니라 모델의 성공률, 속도, 일반화, 안전성을 동시에 결정하는 핵심 설계이다. Continuous regression은 빠르지만 multi-modality에 약하고, action token은 VLM 구조와 잘 맞지만 quantization 문제가 있으며, action chunking은 latency를 줄이지만 반응성을 희생할 수 있다. Diffusion과 flow 기반 정책은 trajectory distribution을 풍부하게 표현하지만 real-time constraint를 해결해야 한다.

다음 장에서는 이런 행동 표현을 가진 VLA를 실제 과제에 맞추는 fine-tuning과 post-training 문제를 다룬다. Action representation이 ``무엇을 출력할 것인가''의 문제라면, fine-tuning은 ``그 출력을 특정 데이터와 과제에 어떻게 맞출 것인가''의 문제이다.

\begin{casebox}{Case study: FAST와 VQ-VLA}
FAST는 per-dimension, per-timestep binning이 high-frequency dexterous data에서 약하다는 문제를 지적하고 DCT 기반 Frequency-space Action Sequence Tokenization을 제안한다. 핵심은 language model에 맞는 discrete token을 쓰되, continuous trajectory의 주파수 구조를 보존하려는 것이다. VQ-VLA는 vector-quantized action tokenizer를 대규모 synthetic/real trajectory로 scaling하여 smoother action과 real-world long-horizon success 개선을 주장한다. 두 논문 모두 ``token이 좋다''가 아니라 ``tokenization이 closed-loop execution을 얼마나 덜 망가뜨리는가''로 읽어야 한다.
\end{casebox}

\begin{table}[t]
  \centering
  \small
  \caption{Claim-source-evidence-caution map for action representation}
  \label{tab:ch4_claim_source}
  \begin{tabularx}{\textwidth}{L{31mm}L{28mm}L{48mm}Y}
    \toprule
    Claim & Source & Evidence & Caution \\
    \midrule
    Action tokenizer는 독립 병목이다 & FAST \citep{fast}, VQ-VLA \citep{vqvla} & DCT tokenization, VQ action tokenizer scaling, closed-loop 또는 long-horizon 결과 & Tokenizer-only effect가 backbone, data, controller 변화와 분리되었는지 확인한다. \\
    Diffusion은 multi-modal action에 강하다 & Diffusion Policy \citep{diffusionpolicy} & Visuomotor trajectory denoising과 real-robot manipulation evidence & VLA backbone 효과가 아니라 action decoder 효과라는 점을 분리한다. \\
    Chunking은 long-horizon imitation에 유용하다 & ACT/ALOHA \citep{aloha} & Bimanual manipulation에서 action chunk prediction과 hardware demonstration & Language grounding이 약한 predecessor이므로 VLA 대표 논문으로 과대분류하지 않는다. \\
    Reconstruction error만으로는 부족하다 & FAST/VQ-VLA plus task results & Token error와 closed-loop success를 함께 봐야 함 & Low reconstruction error가 contact-rich task success를 보장하지 않는다. \\
    \bottomrule
  \end{tabularx}
\end{table}

\begin{sourcebox}{Key papers}
\begin{itemize}
  \item Chi et al. (2023), Diffusion Policy: multimodal continuous action distribution을 trajectory denoising으로 다루는 핵심 선행 논문이다.
  \item Zhao et al. (2023), ACT/ALOHA: action chunking이 long-horizon bimanual manipulation에서 왜 유용한지 보여 준다.
  \item Pertsch et al. (2025), FAST, arXiv preprint: DCT 기반 action tokenization과 FAST+ universal tokenizer를 제안한다.
  \item Wang et al. (2025), VQ-VLA, ICCV 2025: vector-quantized action tokenizer scaling과 real-world long-horizon 개선을 주장한다.
  \item Kim, Finn, and Liang (2025), OpenVLA-OFT, RSS 2025: continuous action representation, action chunking, L1 objective가 fine-tuning 성능에 주는 영향을 연결한다.
\end{itemize}
\end{sourcebox}

\begin{assignmentbox}{Reading assignment [Intermediate]}
FAST와 VQ-VLA를 비교해 codebook/tokenizer가 줄이는 error가 reconstruction error인지, closed-loop task error인지, latency인지 구분하라.
\end{assignmentbox}
